\section{黎曼体积}
\label{sec:mcalc-volume}

前三节的积分对象是$n$形式，结果带符号，依赖定向。日常的面积、体积却是正数，也不需要先决定“正方向”：M\"obius带作为纸条显然有面积，尽管它不可定向。要把函数$f$在流形上累计成$\int f$，例如算球面上的平均气温，需要的是一个正的体积元。光滑结构本身给不出它，因为坐标的尺度可以任意伸缩；黎曼度量规定了长度与角度，也就规定了小平行体的体积。

\subsection{体积密度}
\label{sec:mcalc-volume-density}

\paragraph{小平行体的体积}
在坐标图中，参数空间的小方格$[x^1,x^1+\mathrm dx^1]\times\cdots$近似映成由$\mathrm dx^1\partial_1,\ldots,\mathrm dx^n\partial_n$张成的小平行体。在欧氏空间中，向量$w_1,\ldots,w_n$张成的平行体体积等于Gram矩阵$[\langle w_a,w_b\rangle]$行列式的平方根；把内积换成度量$g$，坐标切向量的Gram矩阵正是$\bm G=[g_{i,j}]$。于是这个小平行体的体积近似为$\sqrt{\det\bm G}\,\mathrm dx^1\cdots\mathrm dx^n$。

\begin{definition}[黎曼体积密度]
\label{def:mcalc-volume-density}
设$(M,g)$是黎曼流形。在坐标图中，称
\begin{equation}
\mathrm dV_g=\sqrt{\det\bm G(x)}\;\mathrm dx^1\cdots\mathrm dx^n
\label{eq:mcalc-volume-density}
\end{equation}
为$g$的\term{体积密度}（volume density），即对坐标域中的可测集$A$规定$\operatorname{vol}_g(A)=\int_{\varphi(A)}\sqrt{\det\bm G}\,\mathrm dx$。
\end{definition}

关键是检查它与坐标无关。换用坐标$y$，记$\bm A=[\partial x^i/\partial y^a]$，由第~\ref{chap:riemannian-geometry}章度量矩阵的变换规律$\widetilde{\bm G}=\bm A^{\top}\bm G\bm A$，取行列式得$\det\widetilde{\bm G}=(\det\bm A)^2\det\bm G$，开方得
\[
\sqrt{\det\widetilde{\bm G}(y)}=\lvert\det\bm A(y)\rvert\,\sqrt{\det\bm G(x(y))}.
\]
平方根自动产生了绝对值，与多元积分换元公式中的$\lvert\det\bm A\rvert$完全匹配。所以同一集合在两张图中算出的体积相同，而这一步没有用到任何定向。与式\eqref{eq:mcalc-top-form-change}对照，$n$形式的系数乘的是带符号的$\det\bm A$，这正是形式积分需要定向而体积密度不需要的原因。

\begin{example}[常见的体积密度]
\label{ex:mcalc-volume-density-examples}
欧氏平面的极坐标中$\bm G=\operatorname{diag}(1,r^2)$，$\mathrm dV=r\,\mathrm dr\,\mathrm d\theta$。单位球面的极角、方位角坐标中$\bm G=\operatorname{diag}(1,\sin^2\theta)$，$\mathrm dV=\sin\theta\,\mathrm d\theta\,\mathrm d\phi$。对嵌入曲面的诱导度量，$\sqrt{\det\bm J^{\top}\bm J}$是参数方格到曲面小块的面积放大率；当$\bm J=\begin{bmatrix}2&1\\0&1\end{bmatrix}$时，$\bm J^{\top}\bm J=\begin{bmatrix}4&2\\2&2\end{bmatrix}$，行列式为$4$，放大率为$2$，与两列$(2,0)^{\top}$、$(1,1)^{\top}$张成的平行四边形面积一致。
\end{example}

\paragraph{可定向情形的体积形式}
若$M$已经定向，可以把体积密度写成一个$n$形式。在每张正向坐标图中令
\begin{equation}
\omega_g=\sqrt{\det\bm G}\,dx^1\wedge\cdots\wedge dx^n .
\label{eq:mcalc-volume-form}
\end{equation}
正向图之间的$\det\bm A>0$，此时$\lvert\det\bm A\rvert=\det\bm A$，由式\eqref{eq:mcalc-top-form-change}，各图的定义在重叠处一致，拼成整体的处处非零的$n$形式，称为\term{体积形式}（volume form）。它有一个不依赖坐标的刻画：对$T_pM$的任一正向标准正交基$(e_1,\ldots,e_n)$，$\omega_g(e_1,\ldots,e_n)=1$。理由是$\omega_g(e_1,\ldots,e_n)$等于$\sqrt{\det\bm G}$乘$e_a$在坐标基下分量矩阵$\bm C$的行列式，而标准正交意味着$\bm C^{\top}\bm G\bm C=\bm I$，即$\det\bm C=1/\sqrt{\det\bm G}$（正向保证取正根）。

不可定向的流形没有处处非零的$n$形式，因而没有体积形式，但式\eqref{eq:mcalc-volume-density}的体积密度照样存在。M\"obius带取任意黎曼度量都有确定的面积。

\subsection{函数积分}
\label{sec:mcalc-function-integration}

\paragraph{整体积分的定义}
有了每张图中的体积密度，就可以仿照形式积分，用单位分解拼接。设$\{(U_\alpha,\varphi_\alpha)\}$是坐标覆盖，$\{\psi_\alpha\}$是从属的单位分解。对非负可测函数$f:M\to[0,\infty]$，定义
\begin{equation}
\int_Mf\,\mathrm dV_g=\sum_\alpha\int_{\varphi_\alpha(U_\alpha)}(\psi_\alpha f)\bigl(\varphi_\alpha^{-1}(x)\bigr)\,\sqrt{\det\bm G_\alpha(x)}\,\mathrm dx\in[0,\infty].
\label{eq:mcalc-function-integral}
\end{equation}
这里“可测”指在每张图中按坐标可测；由于坐标变换光滑，这一条件与图的选取无关。求和对象非负，即使有可数无穷多项也有意义。对一般的可测函数，若$\int\lvert f\rvert\,\mathrm dV_g<\infty$，称$f$绝对可积，分别积分正部与负部再相减。

\begin{theorem}[函数积分的良定性]
\label{thm:mcalc-function-integral-welldefined}
式\eqref{eq:mcalc-function-integral}的值不依赖坐标覆盖与单位分解的选取。由它确定的集函数$\operatorname{vol}_g(A)=\int_M\mathbf 1_A\,\mathrm dV_g$是$M$上Borel集的测度，并且对非负可测或绝对可积的$f$，式\eqref{eq:mcalc-function-integral}就是$f$关于这个测度的Lebesgue积分。
\end{theorem}

\begin{proof}
无关性的论证与定理~\ref{thm:mcalc-form-integral-welldefined}相同：插入另一组单位分解$\{\rho_\beta\}$，把每项写成$\sum_\beta\psi_\alpha\rho_\beta f$的积分，单张图中积分与坐标无关由上一小节的换元计算保证，而交换两个求和的次序由各项非负保证。对绝对可积的$f$，正部与负部分别处理。测度的可数可加性来自单调收敛定理：若$A=\bigsqcup_mA_m$，则$\mathbf 1_A=\sum_m\mathbf 1_{A_m}$，在每张图中积分与求和可以交换，再对$\alpha$求和，非负级数的求和次序也可以交换。最后一句是测度积分的标准构造，可参考第~\ref{chap:integration}章。
\end{proof}

于是第~\ref{chap:integration}章的单调收敛、控制收敛等定理都可以直接用于$\int_Mf\,\mathrm dV_g$。具体计算时常用只漏掉零测集的单张参数化代替单位分解。例如单位球面的面积为
\[
\operatorname{vol}(\mathbb S^2)=\int_0^{2\pi}\!\!\int_0^\pi\sin\theta\,\mathrm d\theta\,\mathrm d\phi=4\pi ,
\]
所用坐标漏掉了一条经线和两极，它们的体积为零。

\paragraph{改变定向、度量与采样测度}
现在可以比较三种看似相近的改动。第一，在定向流形上，由体积形式的刻画，$\int_Mf\,\omega_g=\int_Mf\,\mathrm dV_g$；反转定向时左端的$\omega_g$变成$-\omega_g$，形式积分随之变号，而密度积分$\int f\,\mathrm dV_g$不变。所以“体积”应理解为密度积分，形式积分是它在选定定向后的有符号表示。

第二，改变度量会改变积分。把$g$换成$c^2g$（$c>0$），$\bm G$乘$c^2$，$\sqrt{\det\bm G}$乘$c^n$，体积整体乘$c^n$；半径$a$的球面面积因此是$4\pi a^2$。更一般的度量变化会在不同位置以不同比例改变体积，同一个函数的积分随之改变。

第三，有限样本给出的平均不一定逼近体积积分。设样本点$X_1,X_2,\ldots$独立同分布，其分布关于$\mathrm dV_g$有密度$q$，即$q\geq0$、$\int q\,\mathrm dV_g=1$。若$\int\lvert f\rvert q\,\mathrm dV_g<\infty$，样本平均$\frac1m\sum_{i=1}^mf(X_i)$的期望是$\int fq\,\mathrm dV_g$，而不是$\int f\,\mathrm dV_g/\operatorname{vol}(M)$。例如在单位球面上把$(\theta,\phi)$在$[0,\pi]\times[0,2\pi]$上均匀抽取，坐标空间中的密度是$1/(2\pi^2)$，换成关于$\mathrm dV=\sin\theta\,\mathrm d\theta\,\mathrm d\phi$的密度则是$q=1/(2\pi^2\sin\theta)$，在两极附近趋于无穷：样本挤向两极，用它们的简单平均估计球面平均气温会偏向极地的数值。几何体积与采样测度是两个不同的对象，相差的正是密度$q$。
